% Auto-generated coverage analysis table
\begin{table}[H]
  \centering
  \scriptsize
  \setlength{\tabcolsep}{2pt}
  \renewcommand{\arraystretch}{1.1}
  \begin{tabular}{lcccccccc}
    \toprule
    DGP & N=25 & N=50 & N=100 & N=200 & N=400 & N=800 & N=1600 & N=3200 \\
    \midrule
    TN & (94.7, 93.6) & (93.8, 92.9) & (93.7, 93.2) & (93.7, 92.1) & (92.6, 92.2) & (95.1, 92.1) & (95.5, 92.2) & (91.2, 92.0) \\
    GS3 & (83.7, 90.9) & (89.9, 91.2) & (90.0, 90.8) & (90.5, 91.3) & (91.2, 91.1) & (93.3, 91.4) & (92.2, 91.4) & (91.7, 91.4) \\
    GA3 & (81.5, 91.9) & (85.2, 92.4) & (87.7, 92.6) & (88.6, 92.9) & (90.9, 92.9) & (94.9, 92.8) & (95.6, 92.6) & (95.2, 92.8) \\
    GS5 & (93.8, 95.4) & (91.2, 94.1) & (88.9, 92.7) & (90.4, 93.7) & (95.3, 96.5) & (95.6, 96.5) & (94.5, 96.6) & (94.7, 96.4) \\
    Step & (90.0, 94.2) & (85.9, 92.2) & (81.9, 90.6) & (80.3, 90.3) & (82.9, 90.0) & (90.0, 90.1) & (93.2, 90.5) & (94.2, 91.0) \\
    Sine & (92.5, 94.3) & (91.0, 94.4) & (88.5, 94.7) & (85.5, 94.1) & (89.5, 92.8) & (94.2, 93.5) & (95.9, 93.5) & (96.2, 93.5) \\
    \bottomrule
  \end{tabular}
  \caption{Coverage probabilities (\%) for 95\% confidence intervals across DGPs and sample sizes.}
  \label{tab:coverage_analysis}
\end{table}